% =====================================================================
\subsection{Fase 3: sistemas de juego}

\subsubsection{Arquitectura}

Los sistemas del juego están repartidos en scripts con una responsabilidad cada uno
(Tabla~\ref{tab:arquitectura}). Los objetos no se conocen entre sí: se comunican mediante métodos
estáticos del HUD (por ejemplo, mostrar un aviso, cambiar el objetivo o abrir un diálogo) y mediante
eventos, como el que anuncia que se recogió un objeto. Así, una terminal no necesita saber cómo se
dibuja un diálogo.

\begin{TablaAPA}{tab:arquitectura}{Scripts principales del juego}
\begin{tabularx}{\linewidth}{P{4.4cm}L}
\toprule
Script & Responsabilidad \\
\midrule
\cod{EstadoPartida} & Datos que pasan de un nivel a otro: personaje elegido, salud, discos, inventario, munición, notas leídas y cuerpos registrados; guardado y carga \\
\cod{NivelNexus} & Temporizador, temperatura, objetivo inicial, Game Over, paso al nivel siguiente y Victoria; calcula la malla de navegación al cargar el nivel \\
\cod{ControladorJugador} & Movimiento con controles tipo tanque mediante un \textit{CharacterController} \\
\cod{VidaJugador} & Las tres oportunidades de evasión, el estado herido y la muerte \\
\cod{InteraccionJugador} e \cod{Interactuable} & Detecta lo que está cerca, muestra la acción disponible y la ejecuta con la tecla de interactuar \\
\cod{GestorCamaras} y \cod{CamaraFija} & Las cámaras fijas y el cambio entre ellas \\
\cod{HUDNexus}, \cod{PantallaInventario} y \cod{PantallaFinal} & Interfaz del juego, inventario y pantallas de Game Over y Victoria \\
\cod{ArmaJugador} e \cod{IKTecnico} & Armas, apuntado y animación por cinemática inversa \\
\cod{HabilidadJugador} & Habilidad de cada personaje \\
\cod{AndroideIA} y \cod{JefeSentinel} & Inteligencia artificial de los enemigos \\
\cod{AudioNexus} & Música por capas y efectos de sonido \\
\cod{Cinematica} y \cod{TransicionAscensor} & Cinemáticas y transición entre niveles \\
\cod{ControlesTactiles} & Controles para pantallas táctiles \\
\bottomrule
\end{tabularx}
\end{TablaAPA}

\subsubsection{Cámaras fijas}

Como en el survival horror clásico, el juego no usa una cámara que sigue al personaje sino cámaras
fijas en las salas, que crean tensión porque el jugador oye a los androides pero no siempre los ve.
Solo una cámara está encendida a la vez, lo que además reduce el costo de renderizado.

\begin{itemize}
  \item \textbf{Ángulos por sala.} El constructor de niveles ubica en cada sala una cámara
        principal y evalúa como candidatas las esquinas y los puntos medios de las paredes largas.
        Después de decorar la sala, agrega uno a uno el candidato que ve más puntos del suelo que
        ninguna otra cámara ve, hasta cubrir la sala completa. Un punto solo cuenta como cubierto si
        la cámara ve al técnico de los pies a la cabeza, sin obstáculos y ocupando al menos el 11\,\%
        del alto de la imagen.
  \item \textbf{Cambio de cámara.} Unas 16 veces por segundo, el gestor comprueba si el técnico
        sigue dentro de la imagen y si algo lo tapa. Mientras siga en cuadro, la cámara no cambia;
        si sale del cuadro, queda tapado más de 0,6~s o se ve más pequeño que el 9\,\% de la imagen,
        cambia al ángulo de la sala que lo encuadra mejor (Figura~\ref{fig:angulos}, izquierda). Que el umbral
        de salida sea menor que el de entrada evita que la imagen salte entre dos cámaras.
  \item \textbf{Planos de interés.} La cámara también guía al jugador: al acercarse a un objeto
        importante (terminales, paneles, notas, máquinas, puertas con cerradura, válvulas, objetos
        que se recogen, cadáveres y el ascensor) corta a un plano cercano en tres cuartos que
        muestra el objeto y al técnico (Figuras~\ref{fig:angulos}, derecha, y~\ref{fig:interes2}). Solo se activa si el técnico mira
        hacia el objeto, si el objeto está en la misma sala y si han pasado al menos 1,2~s desde el
        último corte; cuando un objeto se recoge, su plano deja de usarse.
  \item \textbf{Proporciones de pantalla.} Los encuadres están pensados para 16:9; en pantallas más
        estrechas el gestor abre el ángulo vertical para conservar el mismo campo horizontal. La
        cobertura se verificó en 16:9, 4:3 y en la proporción de un celular (19,5:9).
\end{itemize}

\FiguraDoble{068_camara-alternativa-oeste.jpg}{070_plano-interes-terminal.jpg}{fig:angulos}{Ángulo alternativo de la sala de inicio (izquierda) y plano de interés de la terminal (derecha)}{}

\FiguraDoble{071_plano-interes-maquina.jpg}{072_plano-interes-puerta-roja.jpg}{fig:interes2}{Planos de interés de la máquina expendedora (izquierda) y de la puerta roja (derecha)}{}

\subsubsection{Controles tipo tanque}

El técnico avanza y retrocede con \tecla{W} y \tecla{S} (o el stick hacia arriba y abajo) y gira
sobre sí mismo con \tecla{A} y \tecla{D}; \tecla{Shift} lo hace correr y \tecla{Shift + S} da media
vuelta rápida. La dirección no depende de la cámara: aunque el ángulo cambie, el personaje sigue
caminando hacia donde mira, por lo que este tipo de control encaja con las cámaras fijas. Al
empezar el juego aparece un recordatorio con los controles (Figura~\ref{fig:recordatorio}). Las
acciones de movimiento, correr e interactuar se leen del archivo de acciones del Input System,
que se habilita al activarse el jugador, y funcionan igual con teclado, mando y, en móviles, con los
controles táctiles.

\Figura[0.8\linewidth]{067_aviso-controles-tanque.jpg}{fig:recordatorio}{Recordatorio de los controles al empezar el nivel}{}

\subsubsection{Salud, temporizador y HUD}

La salud son tres \textit{oportunidades de evasión}. Al ser alcanzado por un androide, el técnico
pierde una, recibe un empujón y queda invulnerable 2,5~s; con una sola oportunidad camina herido y
los bordes de la pantalla laten en rojo, y con ninguna muere y aparece el Game Over. La bebida
energética recupera una oportunidad y la barrita de proteína, todas.

El temporizador de refrigeración cuenta hacia atrás, y la temperatura del núcleo sube de su valor
inicial hasta el 100\,\% cuando el tiempo se agota. Los cartuchos de refrigerante son coleccionables
con contador que suman 20~s cada uno.

El HUD es solo texto pequeño con una sombra suave, sin recuadros, para no tapar el escenario
(Figura~\ref{fig:hud}): arriba a la izquierda, el objetivo actual; arriba a la derecha, la
refrigeración, la temperatura, los discos recuperados (x/3), los cartuchos, la habilidad y el arma
con su munición. La salud no está fija en pantalla: aparece unos segundos cada vez que cambia y
siempre se puede consultar en el inventario. La sombra del texto es un material de TextMeshPro con
una copia difuminada detrás de cada letra, que se lee sobre cualquier fondo.

\Figura[0.85\linewidth]{029_hud-minimo-solo-texto.jpg}{fig:hud}{HUD del juego: objetivo, refrigeración, temperatura y contadores}{}

\subsubsection{Diálogos e interacción}

Los mensajes de las terminales, las notas y los avisos del sistema aparecen en un cuadro de diálogo
con estilo de terminal: el texto se escribe letra a letra y el juego se pausa mientras se lee
(Figura~\ref{fig:dialogo}). Se avanza con \tecla{E}, \tecla{Espacio}, \tecla{Enter}, un clic, el
mando o tocando la pantalla. Todo lo que se usa con la tecla de interactuar hereda de la clase
\cod{Interactuable}: al acercarse aparece la acción disponible (por ejemplo, \textquote{[E] Usar
terminal}) y el jugador solo interactúa con objetos a menos de 4~m.

\Figura[0.85\linewidth]{030_dialogo-terminal-mantenimiento.jpg}{fig:dialogo}{Cuadro de diálogo de la terminal de mantenimiento}{}

\subsubsection{Inventario}

El inventario es la \textit{terminal portátil del técnico}. Se abre con \tecla{Tab} o \tecla{I},
pausa el juego y tiene un diseño propio con estética de sistema operativo en tonos cian. A la
izquierda muestra el retrato del técnico, un monitor de signos vitales cuya línea de pulso cambia de
velocidad y color según el estado (estable, herido o crítico) y un resumen de la misión. Tiene tres
páginas (Figura~\ref{fig:inventario}):

\begin{itemize}
  \item \textbf{Objetos:} cuadrícula de doce casillas, vista previa del objeto seleccionado y su
        descripción. Desde aquí se usan los objetos de curación y se equipan las armas.
  \item \textbf{Archivo:} las notas leídas, para releerlas.
  \item \textbf{Mapa:} plano esquemático del nivel con el nombre de cada sala, las puertas por
        color, los puntos de interés y la posición del jugador parpadeando.
\end{itemize}

\FiguraDoble{037_inventario-objetos.jpg}{038_inventario-archivo-notas.jpg}{fig:inventario}{Inventario: página de objetos (izquierda) y archivo de notas (derecha)}{}

Los retratos se obtuvieron fotografiando los avatares 3D desde una herramienta de editor que coloca
una cámara frente al hueso de la cabeza en una escena con luz de estudio (Figura~\ref{fig:retratos});
los íconos de los objetos, fotografiando cada prefab con fondo transparente; y el mapa se dibuja a
partir del mismo mapa de texto del nivel. Todos los Canvas usan un escalado que garantiza que la
resolución de referencia de 1920 × 1080 quepa completa en cualquier pantalla.

\FiguraDoble{040_retrato-neyder-renderizado.jpg}{041_retrato-deisy-renderizado.jpg}{fig:retratos}{Retratos de Neyder (izquierda) y Deisy (derecha) generados desde sus avatares}{}

\subsubsection{Guardado de la partida}

La partida se guarda en las terminales de respaldo, que consumen una memoria USB del inventario.
Se guardan el nivel, la posición, el tiempo, la salud, el inventario, la munición, las notas leídas
y los cuerpos ya registrados, mediante el sistema \textit{PlayerPrefs} de Unity, que funciona igual
en Windows, en el navegador y en Android. Cuando hay una partida guardada, el menú muestra la opción
\textit{Continuar}, y la pantalla de Game Over permite cargarla.

\subsubsection{Habilidades de los personajes}

\begin{itemize}
  \item \textbf{Señuelo de red (Neyder):} el repetidor vuela en arco hasta donde cae y, mientras
        emite, muestra un anillo de pulso cian y llama cada medio segundo a los androides cercanos,
        que van a revisarlo en estado de búsqueda. Puede distraer a varios a la vez, pero si un
        androide ve al técnico, lo persigue igual (Figura~\ref{fig:senuelo}).
  \item \textbf{KILLSWITCH.exe (Deisy):} un haz de datos cian salta en zigzag hasta el androide más
        cercano a la vista, que se desploma con el visor y las juntas apagados; al final parpadea,
        se reinicia y vuelve a su ruta (Figura~\ref{fig:killswitch}). Es seguro, pero afecta a uno
        solo y obliga a acercarse.
\end{itemize}

La ficha del selector de personajes explica la habilidad de cada técnico, y el HUD muestra su estado
o los segundos de recarga que faltan.

\FiguraDoble{058_selector-ficha-con-habilidad.jpg}{059_senuelo-de-red-neyder.jpg}{fig:senuelo}{Ficha del selector con la habilidad (izquierda) y el señuelo de red atrayendo a un androide (derecha)}{}

\FiguraDoble{060_killswitch-deisy-rayo.jpg}{061_killswitch-androide-apagado.jpg}{fig:killswitch}{KILLSWITCH.exe: el haz de datos (izquierda) y el androide apagado (derecha)}{}

\subsubsection{Armas y combate}

El técnico puede defenderse con una bayoneta M9, una pistola Beretta M9A1 y una escopeta Remington
870, modelos de Sketchfab con licencia CC BY 4.0
\parencite{erdogay_m9,do_m9a1,milin_870} (Figura~\ref{fig:armas}). Una herramienta de editor los
importa reduciendo sus texturas a 1024~píxeles, crea los materiales URP, invierte el canal verde de
los mapas de normales que venían en formato DirectX, mide cada modelo para dejarlo a su tamaño real
y le agrega un punto de \textit{boca} para el fogonazo, el láser y la trazadora. El arma se coloca
en el hueso de la mano derecha del avatar, así que funciona igual con los dos personajes.

\Figura[0.8\linewidth]{074_armas-vitrina.jpg}{fig:armas}{Las tres armas y las cajas de munición}{}

\begin{itemize}
  \item \textbf{Apuntado asistido:} al apuntar, el técnico gira solo hacia el androide visible más
        cercano dentro de un cono de 50°, y un láser rojo marca el objetivo
        (Figura~\ref{fig:disparos}).
  \item \textbf{Apuntar es comprometerse:} mientras apunta, el técnico no camina y gira más
        despacio.
  \item \textbf{El ruido importa:} cada disparo atrae a los androides cercanos (18~m la pistola,
        24~m la escopeta); el cuchillo es silencioso pero obliga a acercarse.
  \item \textbf{Munición:} el cargador se guarda con la partida y la reserva está en el inventario.
\end{itemize}

\begin{TablaAPA}{tab:armas}{Características de las armas}
\begin{tabularx}{\linewidth}{P{2.4cm}LP{1.6cm}P{1.6cm}P{1.5cm}P{1.6cm}}
\toprule
Arma & Daño & Alcance & Cadencia & Cargador & Ruido \\
\midrule
Bayoneta & 2,5 por tajo, arco de 60° & 1,9 m & 0,75 s & --- & Silencioso \\
Pistola & 2 & 25 m & 0,32 s & 15 & 18 m \\
Escopeta & 7 perdigones de 1,6 que pierden fuerza con la distancia & 14 m & 1 s & 6 & 24 m \\
\bottomrule
\end{tabularx}
\end{TablaAPA}

\FiguraDoble{077_pistola-disparo.jpg}{079_escopeta-disparo.jpg}{fig:disparos}{Disparo de pistola con láser, trazadora y chispas (izquierda) y disparo de escopeta en abanico (derecha)}{}

Cada androide tiene 12 puntos de vida: unos seis disparos de pistola, cinco tajos o uno o dos
disparos de escopeta de cerca. Al recibir un golpe, su visor destella, se aturde un instante y
persigue al técnico; un androide apagado con el KILLSWITCH de Deisy recibe el doble de daño. Al
quedar sin vida cae, sus luces se apagan, deja de bloquear el paso y aparece el aviso
\textquote{UNIDAD SENTINEL NEUTRALIZADA} (Figura~\ref{fig:neutralizado}).

\FiguraDoble{080_cuchillo-ataque.jpg}{081_androide-neutralizado.jpg}{fig:neutralizado}{Tajo de bayoneta (izquierda) y androide neutralizado (derecha)}{}

\subsubsection{Cadáveres y notas}

Como en los survival horror clásicos, las armas y la munición no flotan en el nivel: se obtienen al
registrar los cuerpos de quienes entraron antes. Al pulsar \tecla{E} sobre un cuerpo, el técnico
describe lo que ve, recibe lo que llevaba la víctima y lee su nota, que queda en el archivo del
inventario (Figura~\ref{fig:registro}). El arma o la munición se ven junto al cuerpo o en su mano y
desaparecen al tomarlas, y un cuerpo registrado no vuelve a dar nada, tampoco al cargar la partida.
Las notas cuentan, en conjunto, lo que pasó en el centro de datos (Tabla~\ref{tab:cuerpos}).

\begin{TablaAPA}{tab:cuerpos}{Cuerpos que se pueden registrar en cada nivel}
\begin{tabularx}{\linewidth}{P{0.9cm}P{4.2cm}P{3.2cm}L}
\toprule
Nivel & Cuerpo & Lleva & Lo que cuenta su nota \\
\midrule
1 & Mendoza, técnico & --- & Sacar el disco despierta a las unidades; la salida es el ascensor de la cafetería \\
1 & J. Ramírez, técnico & Bayoneta M9 & Los SENTINEL los encerraron; Mendoza fue por el disco \\
2 & Sargento Ortiz & Escopeta y 3 cartuchos & Orden de misión: Nexus Corp. lo sabía desde el día 14 \\
2 & Operador Castro & 9 balas & Su última transmisión por radio \\
2 & Operador Vargas & Pistola y 9 balas & Registro de la unidad táctica \\
2 & Operador Ruiz & 2 cartuchos & Al sonar la alarma salen todos \\
2 & Guía del equipo & 8 balas & Dónde están el disco y la salida \\
3 & Ibarra & 2 cartuchos & Los anillos de patrulla y algo grande en una cápsula \\
3 & Teniente Salgado & 3 cartuchos y 6 balas & Las válvulas tienen un orden: está en el manual \\
3 & Operador Medina & 12 balas & No gastar las balas en las patrullas \\
3 & Ingeniero Ríos, jefe de sistemas & Barrita de proteína & CERBERUS 7.1 se pidió para proteger los datos a cualquier costo; el reinicio publica los registros \\
\bottomrule
\end{tabularx}
\end{TablaAPA}

\FiguraDoble{084_cadaver-cafeteria-bayoneta.jpg}{085_registro-obtuviste-bayoneta.jpg}{fig:registro}{El técnico Ramírez en la cafetería con la bayoneta (izquierda) y el registro del cuerpo (derecha)}{}

\Figura[0.8\linewidth]{087_nivel2-nota-vargas.jpg}{fig:nota}{Nota del operador Vargas en el Nivel 2}{}

\subsubsection{Controles táctiles}

Para la versión de Android y el navegador del celular se programaron controles táctiles con el
mismo estilo del HUD: anillos finos y texto (Figura~\ref{fig:tactiles}). Se crean solos cuando el
juego corre en un dispositivo móvil y solo se muestran mientras el jugador tiene el control. A la
izquierda hay un joystick que aparece donde se apoya el pulgar, con los mismos controles tipo
tanque, y los botones CORRER e INVENTARIO; a la derecha, INTERACTUAR, HABILIDAD, ARMA y, según el
arma equipada, ATACAR o DISPARAR, APUNTAR y RECARGAR. Si se dispara sin estar apuntando, el técnico
apunta solo al androide más cercano. Los botones usan las mismas entradas que las pruebas
automáticas, de modo que no se duplicó la lógica del juego.

\Figura[0.85\linewidth]{115_controles-tactiles.jpg}{fig:tactiles}{Controles táctiles con la pistola equipada}{}
